% Exemplo: lista de referências escrita à mão com o ambiente thebibliography (sem arquivo .bib)
% Compilar: latexmk -pdf thebibliography.tex       Testado com: TeX Live 2024 ou mais recente
% Forma histórica: a formatação é manual e todos os itens aparecem, citados ou não.
\documentclass[a4paper]{article}
\usepackage[T1]{fontenc}
\usepackage{lmodern}
\usepackage[brazilian]{babel}

\begin{document}

%%BEGIN:thebibliography
O \TeX\ foi criado por Knuth \cite{knuth1986}. Para gráficos, consulte
\cite{goossens1997}.

\begin{thebibliography}{9}   % o argumento é o rótulo mais largo da lista
  \bibitem{knuth1986} KNUTH, D. E. \textbf{The \TeX book}.
    Reading: Addison-Wesley, 1986.
  \bibitem{goossens1997} GOOSSENS, M.; RAHTZ, S.; MITTELBACH, F.
    \textbf{The \LaTeX\ Graphics Companion}. Reading: Addison-Wesley, 1997.
\end{thebibliography}
%%END:thebibliography

\end{document}
